\chapter{The network airlock}\label{ch:network}

\section{The rule}

\vapor{} has no network access of its own. The one exception is the agent backends, which speak to the
hosted model a person configured. Everything else that comes from outside comes through the
\textbf{siphon} (\texttt{Vapor.Siphon}), under one rule: \emph{an agent never connects on its own; only
the person opens the tap.}

\section{Fetchers}

The person declares fetchers in \texttt{\$VAPOR\_HOME/siphons.json}. A fetcher is any program (the
official Python client of a model hub, \texttt{aws s3 cp}, \texttt{rsync} to a private cluster) with an
argument template, the environment variables it may see, a deadline and a byte cap:

\begin{lstlisting}
{"fetchers": [
  {"name": "hub", "argv": ["python3", "fetch.py", "{ref}", "{out}"],
   "env": ["HF_TOKEN"], "timeout_s": 600, "max_bytes": 20000000000},
  {"name": "range", "argv": ["sh", "-c", "curl -r {range} ...", "{ref}", "{out}"]}
]}
\end{lstlisting}

\begin{table}[h]
\centering\small
\begin{tabularx}{\textwidth}{@{}lX@{}}
\toprule
who & what they can do \\
\midrule
the person, at a terminal & \texttt{vapor siphon run NAME REF}; \texttt{approve ID}; \texttt{headers}; \texttt{preflight} \\
an agent (MCP) & \texttt{siphon\_propose}: queue a request with its reason; nothing is fetched \\
the browser console & nothing: the verb is not offered \\
\bottomrule
\end{tabularx}
\caption{Who opens the tap.}
\end{table}

\section{What a fetch can do, and what comes in}

A fetch runs as an operating-system port with only the declared environment variables, a deadline and
a byte cap; exceeding either kills it. A reference that could be read as an option is refused, and a
reference is always one argument, never a shell word. What lands goes through the format airlocks
(safetensors and GGUF headers parsed and bounded, JSON decoded) and may be pinned by SHA-256. Every fetch
leaves a receipt with the exact argument vector, the digests and the formats found; nothing lands from a
malformed file, a wrong digest or a failing fetcher.

\section{Headers before data}

A remote safetensors file can be admitted from its header alone, with two ranged reads: the first eight
bytes give the header length, the next read gives the header. \texttt{vapor siphon preflight NAME CONFIG
SHARD\ldots} fetches the configuration and every shard's header and asks the model airlock whether the
checkpoint can run, naming what is missing or unexpected, before a byte of weights is fetched.

\section{Streaming weights: the arithmetic}

Decoding reads every active weight once per token, so the time per token is bounded below by active
bytes divided by bandwidth.

\begin{table}[h]
\centering\small
\begin{tabular}{@{}lrrrr@{}}
\toprule
model & bytes per token & RAM 50\,GB/s & NVMe 3\,GB/s & net 20\,MB/s \\
\midrule
15\,GB dense & 15\,GB & 0.3\,s & 5\,s & 12.5\,min \\
2.8\,T hybrid (MXFP4 experts) & $\approx$146\,GB & $\approx$3\,s & $\approx$50\,s & $\approx$2\,h \\
\bottomrule
\end{tabular}
\caption{Streaming weights over a network does not serve conversation. It serves work that reads the
weights once for a whole input: one prefill pass to score, classify or embed a long document.}
\end{table}
